\section{Zadanie 16}
Niech \(h\) będzie mierzalną funkcją nieujemną na \(\Omega\).
Udowodnij, że\\ jeśli \(\int_E h\ d\mu = 0\), to
\(\mu(\{x\in E \mid h(x) \neq 0\}) = 0\) (mówimy wtedy, że \(h\mid_E = 0\)
\(\mu\)-prawie wszędzie)
\subsection*{Rozwiązanie}
Parafrazując: wszystko jest nieujemne i sumuje się do zera, więc
(prawie) wszystko jest zerem.

Zdefinujmy \(g: R_{\geqslant0} \to R_{\geqslant0}\) jako
\[
    g(x) = 
    \begin{cases}
        0 & x=0,\\
        \max\{2^a \mid a\in{\mathbb Z}, 2^a \leqslant x\}, & x>0
    \end{cases}
\]
Teraz zdefiniujmy \(f = g \circ h\). Przyglądając się dostatecznie długo 
mamy, że \(f\) jest mierzalne.

Oczywiście \(f \leqslant h\) (bo \(g\leqslant id\)), oraz 
\(\forall_x f(x) = 0 \Leftrightarrow g(x) = 0\).

Korzystając ze schodkowej definicji \(f\), niech \[S_a = 
\{x\in E\mid h(x) \in [2^a, 2^{a+1})\} = \{x\in E\mid f(x) = 2^a\}\]
oraz niech \(\widetilde S = \{x\mid f(x) = h(x) = 0\}\). Wtedy:
\[
    0 = \int_E h\ d\mu \geqslant \int_E f\ d\mu
    = 0\cdot\mu(\widetilde S) + \sum_{a\in\mathbb{Z}} 2^a\mu(S_a)
    = \sum_{a\in\mathbb{Z}} 2^a\mu(S_a)
\]
{\tiny
(łatwo jest wprowadzić na \(\mathbb{Z}\) dobry porządek podobny do
\(\omega\), więc akceptujemy iterowanie sumy po elementach \(\mathbb Z\))
}

Aby suma liczb nieujemnych była zerowa, musi to być sumą samych zer.
\(\forall_a\ \mu(S_a) = 0\)

Wracając do tezy:
\[
    \mu(\{x\in E \mid h(x) \neq 0\})
    = \mu(\{x\in E \mid f(x) \neq 0\})
    = \mu(\bigcup_{a\in\mathbb Z} S_a)
    = \sum_{a\in\mathbb Z} \mu(S_a)
    = 0
\]
